% SPDX-License-Identifier: MIT
% Copyright (c) 2026 Evgeny Baulin
% Optimization Methods, Seminar 04: theory handout (Russian).
% Generated by checks/04. Accelerated Gradient Methods.ipynb from web/04/js/content-*.js and content/students.json. Do not edit by hand.
% Compile with: latexmk -xelatex -interaction=nonstopmode -halt-on-error Theory_ru.tex
\documentclass[11pt,a4paper]{article}

\usepackage{fontspec}
\setmainfont{cmunrm.otf}[BoldFont=cmunbx.otf, ItalicFont=cmunti.otf, BoldItalicFont=cmunbi.otf]
\setsansfont{cmunss.otf}[BoldFont=cmunsx.otf, ItalicFont=cmunsi.otf, BoldItalicFont=cmunso.otf]
\setmonofont{cmuntt.otf}[BoldFont=cmuntb.otf, ItalicFont=cmunit.otf, BoldItalicFont=cmuntx.otf]
\usepackage{polyglossia}
\setdefaultlanguage{russian}
\setotherlanguage{english}
\newfontfamily\cyrillicfont{cmunrm.otf}[Ligatures=TeX, BoldFont=cmunbx.otf, ItalicFont=cmunti.otf, BoldItalicFont=cmunbi.otf]
\newfontfamily\cyrillicfontsf{cmunss.otf}[Ligatures=TeX, BoldFont=cmunsx.otf, ItalicFont=cmunsi.otf, BoldItalicFont=cmunso.otf]
\newfontfamily\cyrillicfonttt{cmuntt.otf}[BoldFont=cmuntb.otf, ItalicFont=cmunit.otf, BoldItalicFont=cmuntx.otf]

\usepackage[a4paper,left=2.3cm,right=2.3cm,top=2.4cm,bottom=2.4cm,headheight=14pt]{geometry}
\usepackage{amsmath,amssymb}
\usepackage{graphicx}
\usepackage{array,booktabs,longtable}
\usepackage{enumitem}
\usepackage{xcolor}
\usepackage[most]{tcolorbox}
\usepackage{titlesec}
\usepackage{fancyhdr}
\usepackage{needspace}
\usepackage[expansion=false]{microtype}
\usepackage[colorlinks=true,linkcolor=hseblue,urlcolor=hseblue,citecolor=hseblue]{hyperref}

\definecolor{hseblue}{RGB}{16,45,105}
\definecolor{hseaccent}{RGB}{198,40,40}
\definecolor{hsegray}{RGB}{108,114,126}
\definecolor{hselight}{RGB}{240,243,248}
\definecolor{hsewarm}{RGB}{253,244,232}

\hypersetup{pdftitle={Методы оптимизации. Семинар 4. Конспект теории}, pdfauthor={Евгений Баулин}}

\titleformat{\section}{\color{hseblue}\Large\bfseries}{\thesection}{0.6em}{}
\titleformat{\subsection}{\color{hseblue}\large\bfseries}{\thesubsection}{0.6em}{}
\titlespacing*{\section}{0pt}{16pt plus 4pt minus 2pt}{8pt plus 2pt}
\titlespacing*{\subsection}{0pt}{12pt plus 3pt minus 2pt}{6pt plus 2pt}
% a heading keeps a few lines under it on its page: a breakable box right after it would otherwise leave it alone
\pretocmd{\section}{\Needspace{7\baselineskip}}{}{\errmessage{cannot patch section}}
\pretocmd{\subsection}{\Needspace{5\baselineskip}}{}{\errmessage{cannot patch subsection}}

\tcbset{handoutbox/.style={enhanced, breakable, sharp corners, boxrule=0pt, left=7pt, right=7pt, top=5pt, bottom=5pt, before skip=8pt plus 2pt, after skip=8pt plus 2pt}}
\newtcolorbox{examplebox}{handoutbox, colback=hselight, colframe=hselight, before upper={\textcolor{hseblue}{\bfseries Пример.}\ }}
\newtcolorbox{whybox}{handoutbox, colback=white, colframe=hseblue, leftrule=2.5pt, top=3pt, bottom=3pt, before upper={\textcolor{hseaccent}{\bfseries Почему.}\ }}
\newtcolorbox{answerbox}{handoutbox, colback=hsewarm, colframe=hsewarm, before upper={\textcolor{hseaccent}{\bfseries Ответ.}\ }}
\newtcolorbox{conventionbox}{handoutbox, colback=hselight, colframe=hseblue, leftrule=2.5pt, before upper={\textcolor{hseblue}{\bfseries Обозначения.}\par\smallskip}}

\newcommand{\handoutfigure}[3]{\par\medskip\noindent\begin{minipage}{\linewidth}\centering
  \includegraphics[width=#1\linewidth]{figures/#2.pdf}\par\smallskip
  \begin{minipage}{0.94\linewidth}\small\textcolor{hseblue}{\bfseries Рисунок.}\ #3\end{minipage}\end{minipage}\par\medskip}
\newcommand{\problemhead}[3]{\par\bigskip\noindent{\color{hseblue}\large\bfseries #1.\ #2}\par\nopagebreak
  \noindent{\small\itshape\color{hsegray}#3}\par\nopagebreak\smallskip}
\newtcolorbox{definitionbox}[1]{handoutbox, colback=white, colframe=hseblue, leftrule=2.5pt, before upper={\textcolor{hseblue}{\bfseries #1}\ }}
\newtcolorbox{theorembox}[1]{handoutbox, colback=hselight, colframe=hseblue, leftrule=2.5pt, before upper={\textcolor{hseblue}{\bfseries #1}\ }}
\newtcolorbox{proofbox}{handoutbox, colback=white, colframe=hsegray, leftrule=1.2pt, top=3pt, bottom=3pt, before upper={\textcolor{hsegray}{\bfseries\itshape Доказательство.}\ }}
\newtcolorbox{notebox}[1]{handoutbox, colback=hselight, colframe=hselight, before upper={\textcolor{hseblue}{\bfseries #1}\par\smallskip}}
\newtcolorbox{trapbox}{handoutbox, colback=hsewarm, colframe=hseaccent, leftrule=2.5pt, before upper={\textcolor{hseaccent}{\bfseries Ловушка.}\ }}
\newtcolorbox{selfcheckbox}{handoutbox, colback=white, colframe=hseaccent, leftrule=2.5pt, top=3pt, bottom=3pt, before upper={\textcolor{hseaccent}{\bfseries Проверьте себя.}\ }}
% a display wider than the line is set as a centred paragraph that may break after a relation
\newcommand{\handoutdisplay}[1]{\sbox0{$\displaystyle #1$}\ifdim\wd0>\linewidth
  \par{\centering\binoppenalty=10000\relpenalty=0\relax$\displaystyle #1$\par}\else\[#1\]\fi}
\newcommand{\selfcheckanswer}{\par\smallskip\textcolor{hsegray}{\bfseries Ответ.}\ }

\setlist[itemize]{itemsep=3pt,topsep=4pt}
\setlist[enumerate]{itemsep=3pt,topsep=4pt}
\setlength{\parindent}{0pt}
\setlength{\parskip}{5pt plus 1pt}
\setlength{\emergencystretch}{2em}
\setcounter{tocdepth}{2}

\pagestyle{fancy}
\fancyhf{}
\fancyhead[L]{\small\color{hsegray}Методы оптимизации}
\fancyhead[R]{\small\color{hsegray}Семинар 4}
\fancyfoot[L]{\small\color{hsegray}Евгений Баулин}
\fancyfoot[C]{\small\thepage}
\fancyfoot[R]{\small\color{hsegray}Лицензия MIT}
\renewcommand{\headrulewidth}{0pt}
\renewcommand{\footrulewidth}{0pt}
\fancypagestyle{plain}{\pagestyle{fancy}}

\begin{document}

\begin{center}
  {\large\color{hsegray}Методы оптимизации\par}
  \vspace{4pt}
  {\large\bfseries\color{hseaccent}Семинар 4\par}
  \vspace{8pt}
  {\LARGE\bfseries\color{hseblue}Ускоренные градиентные методы\par}
  \vspace{8pt}
  {\Large Конспект теории\par}
  \vspace{12pt}
  НИУ ВШЭ, факультет компьютерных наук\par
  2026–2027 учебный год\par
  \vspace{8pt}
  Лектор курса: Андрей Игнатов\par
  Евгений Баулин\par
  Telegram: \href{https://t.me/tarakan_tuc}{@tarakan\_tuc}\par
\end{center}

\noindent{\color{hseblue}\rule{\linewidth}{0.6pt}}

\tableofcontents

\clearpage


Медленный градиентный метод может страдать от вытянутой геометрии, дорогих полных градиентов или шумных оценок. Конспект следует лекции и компьютерному семинару Андрея Игнатова: инерция, координатный спуск, сопряжённые градиенты и стохастические наименьшие квадраты.

\begin{conventionbox}
Матрицы $Q$ и $A$ вещественные, симметричные и положительно определённые, если не сказано иное. $\mu$ и $L$ — крайние собственные значения, $\kappa=L/\mu$, $g=\nabla f$, $r=b-Ax$. В регрессии используем \textbf{половину} средней квадратичной ошибки: $F(w)=\|Xw-y\|^2/(2N)$; у обычной MSE градиент и гессиан вдвое больше.
\end{conventionbox}

% =====================================================================
\section{Геометрия и сравнение методов}\label{sec:geometry}

\begin{definitionbox}{Определение.}
Для $f(x)=\tfrac{1}{2}x^\top Qx-b^\top x$ минимум равен $x^*=Q^{-1}b$. Ошибка $e=x-x^*$ даёт зазор $\tfrac{1}{2}e^\top Qe$ и градиент $Qe$. Запишем $Q=U\Lambda U^\top$ с $U^\top U=I$; все скалярные соотношения ниже действуют на $z=U^\top e$.
\end{definitionbox}

\begin{theorembox}{Теорема.}
GD с постоянным шагом даёт $z_{k+1,i}=(1-\alpha\lambda_i)z_{k,i}$. Сходимость из любой точки выполняется ровно при $0<\alpha<2/L$. Лучший постоянный шаг для квадратичной функции равен $2/(L+\mu)$, а наихудший множитель ошибки — $(\kappa-1)/(\kappa+1)$.
\end{theorembox}

\begin{proofbox}
Каждая координата — геометрическая прогрессия. Её модуль стремится к нулю тогда и только тогда, когда $|1-\alpha\lambda_i|<1$. Чтобы минимизировать наибольший модуль на $[\mu,L]$, уравняйте крайние значения: $1-\alpha\mu=-(1-\alpha L)$.
\end{proofbox}

\begin{whybox}
При большом $\kappa$ шаг, не допускающий взрыва в крутом направлении, почти не меняет пологое. Ортогональный поворот сохраняет собственные значения и число обусловленности, но уводит собственные векторы от координатных осей.
\end{whybox}

\begin{trapbox}
Итерация не является общей единицей работы. GD использует полный градиент, координатный шаг может требовать лишь одной частной производной, CG — умножения матрицы и вектора, а SGD — батча градиентов объектов. Запись полной потери также читает все объекты. Указывайте эти счётчики отдельно; время работы — дополнительная величина, зависящая от компьютера.
\end{trapbox}

% =====================================================================
\section{Инерция на квадратичной функции}\label{sec:momentum}

\begin{definitionbox}{Определение.}
Тяжёлый шар хранит перемещение:
\handoutdisplay{
v_{k+1}=\beta v_k-\alpha g(x_k),\qquad x_{k+1}=x_k+v_{k+1}.
}
 Положите $v_0=0$, то есть $x_{-1}=x_0$, если задача не задаёт другую предыдущую точку. Исключив $v$, получим $x_{k+1}=x_k-\alpha g(x_k)+\beta(x_k-x_{k-1})$.
\end{definitionbox}

\begin{theorembox}{Теорема.}
Для каждого собственного значения $\lambda$ у тяжёлого шара возникает многочлен
\handoutdisplay{
p_\lambda(r)=r^2-(1+\beta-\alpha\lambda)r+\beta.
}
 При $0\le\beta<1$ метод устойчив на всей положительно определённой квадратичной функции ровно при
\handoutdisplay{
0<\alpha<{2(1+\beta)\over L}.
}
 Устойчивость означает сходимость линейного соотношения при любой начальной паре.
\end{theorembox}

\begin{proofbox}
Для вещественного приведённого многочлена $r^2-ar+\beta$ оба корня лежат внутри единичного круга тогда и только тогда, когда $|\beta|<1$, $1-a+\beta>0$ и $1+a+\beta>0$. Подстановка $a=1+\beta-\alpha\lambda$ даёт $\alpha\lambda>0$ и $\alpha\lambda<2(1+\beta)$. Самое сильное ограничение задаёт наибольшее собственное значение.
\end{proofbox}

\begin{theorembox}{Теорема.}
Параметры Поляка минимизируют наихудший радиус корней по спектральному интервалу положительно определённой квадратичной функции:
\handoutdisplay{
\alpha_*={4\over(\sqrt L+\sqrt\mu)^2},\quad \beta_*=\left({\sqrt L-\sqrt\mu\over\sqrt L+\sqrt\mu}\right)^2,\quad \rho_*={\sqrt\kappa-1\over\sqrt\kappa+1}.
}
 Масштаб сложности равен $O(\sqrt\kappa\log(1/\varepsilon))$, в сравнении с $O(\kappa\log(1/\varepsilon))$ для GD.
\end{theorembox}

\begin{proofbox}
Положите $s=\sqrt{\beta_*}=\rho_*$. На концах спектра $1+\beta_*-\alpha_*\mu=2s$ и $1+\beta_*-\alpha_*L=-2s$. Для внутренних собственных значений дискриминант неположителен, поэтому комплексно сопряжённые корни имеют модуль $s$. На концах они кратные. Моды с кратными корнями имеют вид $(c+dk)(\pm s)^k$; простая оценка $\|e_k\|\le s^k\|e_0\|$ отсюда не следует.
\end{proofbox}

\begin{trapbox}
Эти оптимальные параметры относятся к квадратичным функциям. Нельзя переносить ту же глобальную гарантию на любую гладкую сильно выпуклую нелинейную функцию. Меньший радиус корней также не требует убывания функции на каждом начальном шаге.
\end{trapbox}

\problemhead{A1}{Тяжёлый шар: параметры и первые шаги}{У доски}
\textbf{Условие.} Для $f(x)=\tfrac{1}{2}x^\top Qx$, $Q=\operatorname{diag}(1;\ 9)$, положите $x_{-1}=x_0=(1;\ 1)$ и $x_{k+1}=x_k-\alpha Qx_k+\beta(x_k-x_{k-1})$. Найдите оптимальные параметры Поляка для квадратичной функции и спектральный радиус. Вычислите $x_1,x_2$ и значения функции. Обязана ли функция убывать на каждом шаге?

\textbf{Решение.}
\begingroup\sloppy
\begin{enumerate}[label=\textbf{\arabic*.}, leftmargin=*]
  \item Здесь $\mu=1$, $L=9$, $\kappa=9$. Для положительно определённой квадратичной функции
\handoutdisplay{
\alpha_*={4\over(\sqrt L+\sqrt\mu)^2}=0{,}25,\qquad \beta_*=\left({\sqrt L-\sqrt\mu\over\sqrt L+\sqrt\mu}\right)^2=0{,}25.
}
 Максимальный модуль корня равен $\rho=\sqrt{\beta_*}=0{,}5$.
  \item Вдоль собственного вектора с собственным значением $\lambda$ получаем
\handoutdisplay{
z_{k+1}=(1+\beta-\alpha\lambda)z_k-\beta z_{k-1}.
}
 Характеристический многочлен равен $r^2-(1+\beta-\alpha\lambda)r+\beta$. На концах спектра корни кратные: $+\rho$ и $-\rho$.
  \item На первом шаге инерции нет, поскольку $x_0-x_{-1}=0$. Поэтому $x_1=(0{,}75;\ -1{,}25)$; следующий шаг включает $\beta(x_1-x_0)$ и даёт $x_2=(0{,}5;\ 1)$.
  \item $f(x_0)=5$, $f(x_1)=7{,}3125$, $f(x_2)=4{,}625$. На первом шаге функция растёт, хотя все характеристические корни лежат внутри единичного круга. Радиус корней описывает асимптотическое затухание; кратные корни могут добавлять множитель порядка $k$.
\end{enumerate}
\endgroup
\begin{answerbox}
$\alpha_*=0{,}25$, $\beta_*=0{,}25$, $\rho=0{,}5$. Устойчивость не означает монотонности.
\end{answerbox}

\problemhead{A2}{Граница устойчивости метода с инерцией}{Самостоятельно}
\textbf{Условие.} Для $f(z)=\tfrac{1}{2}\lambda z^2$ возьмите $\lambda=16$ и $\beta=0{,}5$. Выведите строгий интервал устойчивых положительных шагов. Укажите его верхнюю границу и максимальный модуль корня при $\alpha=0{,}2$. Устойчив ли этот шаг? Отдельно разберите равенство на границе.

\textbf{Решение.}
\begingroup\sloppy
\begin{enumerate}[label=\textbf{\arabic*.}, leftmargin=*]
  \item При $a=1+\beta-\alpha\lambda$ критерий для корней сводится к $|\beta|<1$, $\alpha\lambda>0$ и $2(1+\beta)-\alpha\lambda>0$. При заданной неотрицательной инерции $0<\alpha<2(1+\beta)/\lambda=0{,}1875$.
  \item На границе многочлен раскладывается в $(r+1)(r+\beta)$, поэтому один корень равен $-1$. При обычной начальной паре сохраняется колеблющаяся компонента. Для предложенного большего шага максимальный модуль равен $1{,}3217>1$, так что для некоторых начальных точек метод расходится.
\end{enumerate}
\endgroup
\begin{answerbox}
Верхняя граница: $0{,}1875$; радиус для пробного шага: $1{,}3217$. Интервал устойчивости открытый.
\end{answerbox}

\handoutfigure{0.95}{fig_momentum}{Нетривиальный поворот показывает крутое и пологое направления, не совмещая их с осями.}

\handoutfigure{0.9}{fig_stability}{Аналитическая граница строгая; конечный запуск зависит ещё и от порога уменьшения.}

% =====================================================================
\section{Опережение и рестарт}\label{sec:lookahead}

\begin{definitionbox}{Определение.}
Форма с постоянной инерцией и опережением задаётся как
\handoutdisplay{
y_k=x_k+\beta(x_k-x_{k-1}),\qquad x_{k+1}=y_k-\alpha g(y_k).
}
 На скалярной квадратичной функции
\handoutdisplay{
z_{k+1}=(1-\alpha\lambda)\big((1+\beta)z_k-\beta z_{k-1}\big).
}
 Многочлен отличается от тяжёлого шара: $r^2-(1+\beta)(1-\alpha\lambda)r+\beta(1-\alpha\lambda)$.
\end{definitionbox}

\begin{notebox}{Два режима ускорения}
Для известной гладкой сильно выпуклой функции стандартный ускоренный выбор равен $\alpha=1/L$, $\beta=(\sqrt L-\sqrt\mu)/(\sqrt L+\sqrt\mu)$. В просто выпуклом случае ускорение использует меняющийся коэффициент экстраполяции; оценка зазора $O(1/k^2)$ не относится к произвольной постоянной инерции.
\end{notebox}

\begin{whybox}
Если экстраполированный шаг направлен неверно, отбросьте инерцию и повторите градиентный шаг из текущей точки. Рестарт по функции проверяет, превышает ли предложенное значение текущее. Обозначайте метод с рестартом отдельно и учитывайте дополнительное вычисление градиента.
\end{whybox}

\problemhead{A3}{Где вычисляется градиент?}{Самостоятельно}
\textbf{Условие.} Для $f(x)=\tfrac{1}{2}\lambda x^2$ положите $\lambda=4$, $x_k=2$, $x_{k-1}=3$, $\alpha=0{,}1$, $\beta=0{,}5$. Найдите следующую точку тяжёлого шара и следующую точку метода с опережением: $y_k=x_k+\beta(x_k-x_{k-1})$, $x_{k+1}=y_k-\alpha\nabla f(y_k)$. Объясните различие.

\textbf{Решение.}
\begingroup\sloppy
\begin{enumerate}[label=\textbf{\arabic*.}, leftmargin=*]
  \item Экстраполированная точка равна $y_k=1{,}5$. Тяжёлый шар использует $\nabla f(x_k)=8$ и даёт $x_k+\beta(x_k-x_{k-1})-\alpha\nabla f(x_k)=0{,}7$.
  \item Обновление с опережением берёт $\nabla f(y_k)=6$ и даёт $y_k-\alpha\nabla f(y_k)=0{,}9$. Если в обеих реализациях брать градиент в $x_k$, методы незаметно станут одинаковыми.
\end{enumerate}
\endgroup
\begin{answerbox}
Тяжёлый шар: $0{,}7$. Опережение: $0{,}9$.
\end{answerbox}

% =====================================================================
\section{Координатный спуск}\label{sec:coordinates}

\begin{definitionbox}{Определение.}
Координатная константа Липшица $L_i>0$ даёт
\handoutdisplay{
f(x+te_i)\le f(x)+tg_i+\tfrac{1}{2}L_it^2.
}
 Выбор $t=-g_i/L_i$ приводит к
\handoutdisplay{
f(x^+)\le f(x)-{g_i^2\over2L_i}.
}
 Для квадратичной функции $L_i=Q_{ii}$, и неравенство становится равенством.
\end{definitionbox}

\begin{proofbox}
Для квадратичной функции срез раскрывается точно: $f(x+te_i)=f(x)+tg_i+Q_{ii}t^2/2$. Дифференцирование по $t$ даёт единственный минимум. Подстановка даёт уменьшение. После изменения $x_i$ следующая производная должна использовать изменённый вектор.
\end{proofbox}

\begin{notebox}{Три правила выбора}
Циклический выбор проходит координаты по порядку. Равномерный случайный выбор выбирает одну координату на шаг. Жадный выбор максимизирует $|g_i|/\sqrt{L_i}$, поскольку его квадрат равен удвоенному гарантированному уменьшению. В размерности $d$ один эквивалентный проход — это $d$ координатных обновлений, но проход не обязательно стоит столько же, сколько один полный градиент.
\end{notebox}

\begin{theorembox}{Теорема.}
Если $f$ является $\mu$-сильно выпуклой и для всех координат $L_i\le L_{\max}$, равномерный выбор даёт
\handoutdisplay{
\mathbb E[f(x^+)-f^*\mid x]\le\left(1-{\mu\over dL_{\max}}\right)(f(x)-f^*).
}
 Это оценка на одно координатное обновление при новой независимой выборке координаты.
\end{theorembox}

\begin{proofbox}
Усредните $g_i^2/(2L_i)$ по координатам: среднее уменьшение не меньше $\|g\|^2/(2dL_{\max})$. Сильная выпуклость даёт $\|g\|^2\ge2\mu(f-f^*)$; подставьте это в оценку уменьшения.
\end{proofbox}

\problemhead{B1}{По одной координате}{У доски}
\textbf{Условие.} Для $f(x)=\tfrac{1}{2}x^\top Qx$, $Q=\begin{pmatrix} 5 & -4 \\ -4 & 5 \end{pmatrix}$, начните из $x_0=(1;\ 1)$. Минимизируйте сначала по координате $1$, затем по координате $2$, фиксируя другую. Выведите точный координатный шаг и найдите уменьшение функции после каждого обновления.

\textbf{Решение.}
\begingroup\sloppy
\begin{enumerate}[label=\textbf{\arabic*.}, leftmargin=*]
  \item Вдоль координаты $i$ имеем $f(x+te_i)=f(x)+tg_i+\tfrac{1}{2}Q_{ii}t^2$. Минимум достигается при $t=-g_i/Q_{ii}$, а уменьшение равно $g_i^2/(2Q_{ii})$. Здесь $L_i=Q_{ii}$, хотя собственные значения равны $(1;\ 9)$.
  \item $g(x_0)=(1;\ 1)$. После координаты $1$: $x=(0{,}8;\ 1)$, $f=0{,}9$, уменьшение $0{,}1$. Перед следующим шагом пересчитайте градиент в этой точке.
  \item После координаты $2$: $x=(0{,}8;\ 0{,}64)$, $f=0{,}576$, уменьшение $0{,}324$. Это два координатных обновления, то есть один полный проход в двумерной задаче.
\end{enumerate}
\endgroup
\begin{answerbox}
$x_1=(0{,}8;\ 1)$, $x_2=(0{,}8;\ 0{,}64)$. Каждое обновление минимизирует текущий одномерный срез.
\end{answerbox}

\problemhead{B2}{Выбор координаты по уменьшению функции}{Самостоятельно}
\textbf{Условие.} Для $Q=\operatorname{diag}(1;\ 9)$ и $x=(2;\ 1)$ вычислите обе компоненты градиента. Какая координата даёт наибольшее точное уменьшение $g_i^2/(2Q_{ii})$? Укажите номер координаты, уменьшение и новое значение функции. Объясните показатель $|g_i|/\sqrt{Q_{ii}}$.

\textbf{Решение.}
\begingroup\sloppy
\begin{enumerate}[label=\textbf{\arabic*.}, leftmargin=*]
  \item Градиент равен $(2;\ 9)$. Нормированные показатели равны $(2;\ 3)$, поэтому координата $2$ даёт наибольшее точное улучшение.
  \item Обновление этой координаты уменьшает $f$ на $4{,}5$ и оставляет $f=2$. Максимизация $|g_i|/\sqrt{Q_{ii}}$ равносильна максимизации его квадрата, а значит, и точного координатного уменьшения.
\end{enumerate}
\endgroup
\begin{answerbox}
Координата $2$; уменьшение $4{,}5$; новое значение $2$.
\end{answerbox}

% =====================================================================
\section{Сопряжённые градиенты}\label{sec:cg}

\begin{definitionbox}{Определение.}
При $A\succ0$ минимизация $f(x)=\tfrac{1}{2}x^\top Ax-b^\top x$ равносильна решению $Ax=b$. Направления $p_i,p_j$ являются $A$-сопряжёнными, если $p_i^\top Ap_j=0$. Невязка $r=b-Ax$ равна отрицательному градиенту.
\end{definitionbox}

\begin{definitionbox}{Определение.}
Начните с $r_0=b-Ax_0$, $p_0=r_0$. На каждом шаге вычислите
\handoutdisplay{
\begin{gathered}
q_k=Ap_k,\quad \alpha_k={r_k^\top r_k\over p_k^\top q_k}, \\[6pt]
x_{k+1}=x_k+\alpha_kp_k,\quad r_{k+1}=r_k-\alpha_kq_k, \\[6pt]
\beta_k={r_{k+1}^\top r_{k+1}\over r_k^\top r_k},\quad p_{k+1}=r_{k+1}+\beta_kp_k.
\end{gathered}
}
 Если новая невязка нулевая, проверьте завершение до вычисления следующего коэффициента.
\end{definitionbox}

\begin{theorembox}{Теорема.}
В точной арифметике разные невязки ортогональны, а направления поиска являются $A$-сопряжёнными. CG минимизирует квадратичную функцию на $x_0+\mathcal K_k(A,r_0)$, где $\mathcal K_k=\operatorname{span}\{r_0,Ar_0,\ldots,A^{k-1}r_0\}$. В размерности $d$ он завершается не более чем за $d$ шагов.
\end{theorembox}

\begin{whybox}
Минимизация вдоль нового сопряжённого направления не отменяет минимизации вдоль предыдущих: $p_i^\top A(x+tp_j)=p_i^\top Ax$ при $i\ne j$. Когда направления поиска порождают всё пространство, градиент ортогонален всему пространству и должен быть нулевым.
\end{whybox}

\begin{theorembox}{Теорема.}
Для энергетической нормы $\|e\|_A=\sqrt{e^\top Ae}$
\handoutdisplay{
\|e_k\|_A\le2\left({\sqrt\kappa-1\over\sqrt\kappa+1}\right)^k\|e_0\|_A.
}
 Так как $f(x_k)-f^*=\|e_k\|_A^2/2$, в отношении зазоров множитель возводится в квадрат. Группировка собственных значений может ускорить наблюдаемую сходимость по сравнению с наихудшей оценкой по интервалу.
\end{theorembox}

\begin{trapbox}
Используйте $\|r\|/\|b\|$ при $b\ne0$. При $b=0$ задайте явный абсолютный порог. На небольших плотных примерах проверьте симметрию и положительную определённость; предусмотрите проверки неположительной кривизны, неконечных значений и лимита итераций. Истинная конечная невязка равна $b-Ax$, и её проверка стоит ещё одного произведения матрицы на вектор.
\end{trapbox}

\problemhead{B3}{Два шага сопряжённых градиентов}{Самостоятельно}
\textbf{Условие.} Решите $Ax=b$ методом CG при $A=\begin{pmatrix} 4 & 1 \\ 1 & 3 \end{pmatrix}$, $b=(1;\ 2)$, $x_0=(0;\ 0)$. Начните с $r_0=b-Ax_0$, $p_0=r_0$. Вычислите $\alpha_0$, $\beta_0$ и $x_2$. Проверьте $r_0^\top r_1=0$ и $p_0^\top Ap_1=0$. Сравните с прямым решением.

\textbf{Решение.}
\begingroup\sloppy
\begin{enumerate}[label=\textbf{\arabic*.}, leftmargin=*]
  \item Первая минимизация вдоль направления даёт $\alpha_0=0{,}25$, $x_1=(0{,}25;\ 0{,}5)$, $r_1=(-0{,}5;\ 0{,}25)$. Скалярное произведение с $r_0$ равно $0$.
  \item $\beta_0=0{,}0625$, $p_1=r_1+\beta_0p_0=(-0{,}4375;\ 0{,}375)$. Направления удовлетворяют $p_0^\top Ap_1=0$. Следующий шаг равен $\alpha_1=0{,}3636$.
  \item $x_2=x^*=(0{,}0909;\ 0{,}6364)$ и $r_2=0$ в точной арифметике. Для положительно определённой матрицы CG завершается не более чем за размерность шагов, но округление может потребовать дополнительных шагов. Вычисленное произведение $Ap_k$ нужно использовать и в знаменателе, и при обновлении невязки.
\end{enumerate}
\endgroup
\begin{answerbox}
$\alpha_0=0{,}25$, $\beta_0=0{,}0625$, $x_2=(0{,}0909;\ 0{,}6364)$.
\end{answerbox}

\handoutfigure{0.95}{fig_directions}{GD, циклический координатный спуск и CG на одной положительно определённой квадратичной функции с разными единицами работы.}

% =====================================================================
\section{Мини-батчи и шумовой порог}\label{sec:sgd}

\begin{definitionbox}{Определение.}
Для $F(w)=N^{-1}\sum_i\ell_i(w)$ новый равномерный индекс даёт $\mathbb E[\nabla\ell_i(w)\mid w]=\nabla F(w)$. Градиент батча — среднее градиентов его объектов. Если $B$ выборок независимы и каждая имеет дисперсию $\sigma^2$, дисперсия батча равна $\sigma^2/B$.
\end{definitionbox}

\begin{proofbox}
Среднее выборки линейно, поэтому его ожидание равно полному среднему. Раскройте квадрат центрированного среднего: независимые смешанные члены имеют нулевое ожидание, остаются $B$ диагональных членов, делённых на $B^2$. Без возвращения ковариационные члены отрицательны и дают поправку конечной совокупности $(N-B)/(N-1)$.
\end{proofbox}

\begin{notebox}{Правило выборки}
Перемешивание без возвращения в каждой эпохе удобно для реализации. Однако итерат зависит от предыдущих батчей, поэтому условная несмещённость каждого последующего батча не возникает автоматически. Указывайте, предполагает ли теоретический вывод новые независимые батчи или случайное перемешивание.
\end{notebox}

\begin{theorembox}{Теорема.}
Для $L$-гладкой функции и несмещённого градиента с условной дисперсией шума не выше $\sigma^2/B$ лемма о спуске даёт
\handoutdisplay{
\mathbb E[F(w^+)\mid w]\le F(w)-\alpha(1-L\alpha/2)\|\nabla F(w)\|^2+{L\alpha^2\sigma^2\over2B}.
}
 Положительный шумовой член объясняет, почему постоянный шаг может оставлять порог.
\end{theorembox}

\begin{proofbox}
Подставьте $w^+=w-\alpha g$ в верхнюю модель гладкости. Используйте $\mathbb E[g\mid w]=\nabla F(w)$ и $\mathbb E[\|g\|^2\mid w]\le\|\nabla F(w)\|^2+\sigma^2/B$. Соберите градиентные и дисперсионные члены.
\end{proofbox}

\problemhead{C1}{Батч оценивает градиент}{У доски}
\textbf{Условие.} Пусть $\ell_i(w)=\tfrac{1}{2}(w-a_i)^2$, $a=(0;\ 2;\ 4;\ 6)$, а $F(w)$ — средняя функция потерь. В точке $w_0=0$ сравните полный шаг GD с шагом по батчу с целями $(0;\ 2)$, используя $\alpha=0{,}25$. Найдите дисперсию одного равномерно выбранного градиента. Как она изменяется при усреднении двух выборок с возвращением и без?

\textbf{Решение.}
\begingroup\sloppy
\begin{enumerate}[label=\textbf{\arabic*.}, leftmargin=*]
  \item В начальной точке полный градиент равен $-3$, а градиент выбранного батча — $-1$. Поэтому $w_{\mathrm{full}}=0{,}75$ и $w_B=0{,}25$. Батч даёт оценку, а не точный полный градиент.
  \item Минимум достигается на среднем целевом значении $w^*=3$. Вычтите полный градиент из каждого градиента объекта и возведите в квадрат; среднее равно $\sigma^2=5$. Усреднение двух независимых выборок даёт $2{,}5$.
  \item Для равномерного подмножества размера $B$ без возвращения $\operatorname{Var}(g_B)=\sigma^2(N-B)/(B(N-1))$, поэтому здесь дисперсия равна $1{,}6667$. Внутри перемешанной эпохи последующие батчи зависят от предыдущих выборок; условную независимость нельзя предполагать автоматически.
\end{enumerate}
\endgroup
\begin{answerbox}
$g_B=-1$, $w_B=0{,}25$, $\sigma^2=5$.
\end{answerbox}

\problemhead{C2}{Шумовой порог при постоянном шаге}{Самостоятельно}
\textbf{Условие.} Для $F(w)=\tfrac{1}{2}\mu w^2$ градиент батча равен $\mu w+\xi_k$, где шумы независимы, имеют нулевое среднее и дисперсию $\sigma^2/B$. Пусть $\mu=2$, $\sigma^2=9$, $B=3$, $\alpha=0{,}1$. Выведите рекуррентное соотношение для $\mathbb E w_k^2$ и его стационарное значение. Укажите стационарную дисперсию и средний зазор функции.

\textbf{Решение.}
\begingroup\sloppy
\begin{enumerate}[label=\textbf{\arabic*.}, leftmargin=*]
  \item $w_{k+1}=(1-\alpha\mu)w_k-\alpha\xi_k$. Независимость и нулевое среднее дают
\handoutdisplay{
\mathbb E w_{k+1}^2=(1-\alpha\mu)^2\mathbb E w_k^2+\alpha^2\sigma^2/B.
}
 Здесь среднее сжимается с множителем $0{,}8$.
  \item При $0<\alpha\mu<2$ решите уравнение неподвижной точки:
\handoutdisplay{
V_\infty={\alpha\sigma^2\over B\mu(2-\alpha\mu)}=0{,}0833.
}
 Предельное среднее равно нулю, поэтому этот второй момент — дисперсия; $\mathbb E[F(w)-F^*]=\mu V_\infty/2=0{,}0833$.
  \item Меньший шаг или больший батч понижает порог. Среднеквадратичный радиус равен $0{,}2887$, а не самой дисперсии. Убывающий шаг при подходящих предположениях может убрать постоянный порог.
\end{enumerate}
\endgroup
\begin{answerbox}
$V_\infty=0{,}0833$, средний зазор $0{,}0833$.
\end{answerbox}

% =====================================================================
\section{Расписания и воспроизводимые эксперименты}\label{sec:schedules}

\begin{definitionbox}{Определение.}
Определяйте расписание по \textbf{числу обновлений} $k$, начиная с нуля: постоянное $\alpha_k=\alpha_0$, шаг, обратно пропорциональный корню, $\alpha_k=\alpha_0/\sqrt{1+k}$, и обратное расписание $\alpha_k=\alpha_0/(1+k)$. Счётчик эпох задаёт другое расписание; не смешивайте соглашения.
\end{definitionbox}

\begin{notebox}{Что гарантирует расписание}
При обычных предположениях стохастической аппроксимации достаточные условия на шаг имеют вид $\sum_k\alpha_k=\infty$ и $\sum_k\alpha_k^2<\infty$. Обратные шаги выполняют оба условия; обратный квадратный корень не даёт суммируемости квадратов. Это не запрещает его в выпуклых экспериментах с конечным бюджетом, но не позволяет ссылаться на ту же теорему без дополнительных аргументов.
\end{notebox}

\begin{trapbox}
Обратное расписание может слишком быстро уменьшить шаг ещё до прохождения вытянутой долины. Сравнивайте одинаковые начальный шаг, счётчик обновлений, данные, seed и бюджет работы. Расписание, выигравшее на коротком запуске, может проиграть на длинном.
\end{trapbox}

\begin{notebox}{Протокол эпохи}
В каждой эпохе создавайте перестановку, берите срезы размера $B$ и сохраняйте последний срез, даже если он короче. Используйте фактическую длину при усреднении градиента. Записывайте полную функцию после окончания эпохи. Повторите запуск на $5$ случайных seed и покажите медиану и межквартильный интервал.
\end{notebox}

\handoutfigure{0.9}{fig_sgd}{Используйте одинаковый бюджет градиентов объектов для всех расписаний; полосы показывают разброс по seed.}

% =====================================================================
\section{Стандартизация признаков}\label{sec:scaling}

\begin{definitionbox}{Определение.}
Для центрированной матрицы $X$ пусть $S=\operatorname{diag}(s_i)$ содержит положительные стандартные отклонения обучающих признаков, а $Z=XS^{-1}$. Гессиан задачи наименьших квадратов равен $H_Z=S^{-1}HS^{-1}$ при $H=X^\top X/N$. Для равенства предсказаний $Zu=Xw$ нужно $w=S^{-1}u$.
\end{definitionbox}

\begin{whybox}
Столбцы с разными единицами вытягивают квадратичную функцию. Масштабирование может убрать эту часть числа обусловленности. Оно не устраняет почти линейную зависимость столбцов, поэтому стандартизированный гессиан может остаться плохо обусловленным. Для постоянных столбцов нужно отдельное соглашение вместо деления на нуль.
\end{whybox}

\begin{notebox}{Вернуть исходные единицы}
При средних $m$ стандартизированное предсказание $c+(x-m)^\top S^{-1}u$ превращается в $b+x^\top w$, где $w=S^{-1}u$ и $b=c-m^\top w$. Оценивайте $m,S$ только по обучающим данным. Поиск наибольшего пригодного шага SGD по сетке — эмпирическая проверка конечного запуска, отличающаяся от границы полного градиента $2/\lambda_{\max}(H)$.
\end{notebox}

\problemhead{C3}{Масштаб признаков и возврат единиц}{Самостоятельно}
\textbf{Условие.} Для центрированной задачи наименьших квадратов $F(w)=\|Xw-y\|^2/(2N)$ пусть $H=X^\top X/N=\begin{pmatrix} 1 & 0 \\ 0 & 100 \end{pmatrix}$. Стандартизуйте с помощью $Z=XS^{-1}$, $S=\operatorname{diag}(1;\ 10)$. Найдите числа обусловленности до и после масштабирования и границы устойчивости полного GD. Если в стандартизированных единицах получено $u=(2;\ 30)$, чему равны коэффициенты в исходных единицах?

\textbf{Решение.}
\begingroup\sloppy
\begin{enumerate}[label=\textbf{\arabic*.}, leftmargin=*]
  \item $H_Z=S^{-1}HS^{-1}=I$. Поэтому $\kappa(H)=100$, $\kappa(H_Z)=1$. Верхние границы шага полного GD равны соответственно $2/L=0{,}02$ и $2$; обе границы исключаются.
  \item Верните коэффициенты: $w=S^{-1}u=(2;\ 3)$. Для нецентрированных признаков со средним $m$ стандартизированная константа $c$ превращается в $c-m^\top w$. Совпадения коэффициентов в двух системах координат ожидать не нужно.
  \item В этом диагональном примере масштабирование исправляет разные единицы. Сильные корреляции могут остаться после стандартизации, а граница по полному гессиану не гарантирует устойчивость каждого отдельного случайного батча.
\end{enumerate}
\endgroup
\begin{answerbox}
$\kappa$: $100\to1$; $w=(2;\ 3)$.
\end{answerbox}

\handoutfigure{0.9}{fig_scaling}{Потери в исходных и стандартизированных единицах при одинаковых настройках SGD.}

% =====================================================================
\section{Проверки реализации и упражнения}\label{sec:implementation}

\begin{itemize}
  \item Возвращайте независимую копию каждой сохранённой точки; добавление одного изменяемого массива стирает траекторию.
  \item Останавливайтесь по текущему градиенту или невязке, проверяйте конечность значений и указывайте лимит итераций.
  \item Указывайте отдельно работу инициализации, оптимизации и диагностики. Жадный просмотр координат вычисляет больше, чем одну выбранную производную.
  \item Сравните реализации NumPy с точными маленькими примерами и прямой линейной алгеброй, затем запускайте большие эксперименты.
\end{itemize}

\begin{selfcheckbox}
Почему аналитически устойчивая пара может не пройти проверку скорости за конечное число шагов?
\selfcheckanswer Модуль корня возле единицы и начальный переходный режим могут удерживать функцию выше заданного порога конечного запуска.
\end{selfcheckbox}

\begin{selfcheckbox}
Почему одну итерацию CG нельзя напрямую сравнить с одним координатным обновлением?
\selfcheckanswer CG вычисляет произведение матрицы на вектор, тогда как координатный шаг может использовать одну частную производную. Реальная стоимость зависит от разреженности и реализации выбора координаты.
\end{selfcheckbox}

\end{document}
